%% =================================================================
%% Beber et al. IEEE Robotics and Automation Letters, Feb 2026.
%% Reconstruction of page 7 (printed p. 7).
%% Contents: Table II placeholder (100-run performance: DR, time, RMSE
%% across BO-EI / BO-EIS / Ergodic, trajectory lengths 500/650/800 mm,
%% w/ and w/o noise); Table III placeholder (segmentation
%% Sens/Spec for three clusters); Fig. 4 placeholder (clustering +
%% boundary extraction results); Section IV-D continuation (50-run
%% simulation + sensitivity/specificity defn); Section V Experimental
%% Results (instrumentation rollcall, silicone-phantom 5-experiment
%% validation results).
%% =================================================================

\documentclass[11pt]{article}

\usepackage[margin=0.85in]{geometry}
\usepackage{amsmath, amssymb}
\usepackage{mathrsfs}
\usepackage{xcolor}
\usepackage{fancyhdr}

\pagestyle{fancy}
\fancyhf{}
\lhead{\small Beber et al.: Ergodic Palpation for Tissue Mapping}
\rhead{\small 7}
\cfoot{\small 7 $\to$ \arabic{page}}
\renewcommand{\headrulewidth}{0.4pt}

\setcounter{page}{1}
\setcounter{equation}{5}
\setlength{\parskip}{6pt}
\setlength{\parindent}{0pt}

\begin{document}

% ---------- Table II ----------
\begin{center}
\textbf{TABLE II}\\[2pt]
\textsc{Comparison of performance metrics across 100 runs of Ergodic
Search and BO, with (w/) and without (w/o) noise. DR is the Detection
Rate, $t$ the Exploration Time, and RMSE the Root Mean Square Error.}
\end{center}

\begin{center}
\footnotesize
\begin{tabular}{|l|ccc|ccc|ccc|}
\hline
Traj.\ Length
  & \multicolumn{3}{c|}{500 mm}
  & \multicolumn{3}{c|}{650 mm}
  & \multicolumn{3}{c|}{800 mm} \\
\hline
Method & DR $\uparrow$ & $t\downarrow$ & RMSE$\downarrow$
       & DR $\uparrow$ & $t\downarrow$ & RMSE$\downarrow$
       & DR $\uparrow$ & $t\downarrow$ & RMSE$\downarrow$ \\
\hline
BO-EI (w/o)   & 39 & 59 & 12.09 & 73 & 88 & 7.98 & 92 & 124 & 5.47 \\
BO-EIS (w/o)  & 93 & 60 &  4.61 & 99 & 92 & 2.81 & 99 & 124 & 1.87 \\
Ergodic (w/o) & 78 & 41 &  4.69 & 92 & 62 & 2.74 & 99 &  84 & 1.42 \\
BO-EI (w/)    & 38 & 60 & 12.10 & 60 & 89 & 7.84 & 91 & 118 & 5.49 \\
BO-EIS (w/)   & 85 & 63 &  5.54 & 95 & 95 & 3.35 & 99 & 130 & 2.59 \\
Ergodic (w/)  & 79 & 41 &  5.87 & 98 & 62 & 3.03 & 99 &  84 & 2.16 \\
\hline
\end{tabular}
\end{center}

\vspace{0.6em}

% ---------- Table III ----------
\begin{center}
\textbf{TABLE III}\\[2pt]
\textsc{Average segmentation results across three clusters.}
\end{center}

\begin{center}
\footnotesize
\begin{tabular}{|c|ccc|ccc|ccc|}
\hline
 & \multicolumn{3}{c|}{Cluster 1}
 & \multicolumn{3}{c|}{Cluster 2}
 & \multicolumn{3}{c|}{Cluster 3} \\
\hline
 & BO-EI & BO-EIS & Erg. & BO-EI & BO-EIS & Erg. & BO-EI & BO-EIS & Erg. \\
\hline
Sens. & 0.561 & 0.950 & 0.975 & 0.657 & 0.997 & 0.971 & 0.746 & 0.904 & 0.978 \\
Spec. & 0.654 & 0.993 & 0.993 & 0.655 & 0.998 & 0.997 & 0.829 & 0.953 & 0.992 \\
\hline
\end{tabular}
\end{center}

\vspace{0.8em}

% ---------- Fig. 4 placeholder ----------
\begin{center}
\fbox{\begin{minipage}[c][4.5cm][c]{0.80\textwidth}
\centering
\textbf{Fig. 4.}\ Clustering and boundary extraction after exploration.\\[4pt]
\textcolor{gray}{\footnotesize
Original cluster shapes (in blue) overlaid against estimated shapes
(in red) for the three segmentation benchmarks. Cluster 1: irregular
inclusion; Cluster 2: circular inclusion; Cluster 3: donut-shaped
inclusion. The overlay visualises boundary fidelity relative to
ground truth.}
\end{minipage}}
\end{center}

\vspace{0.4em}

\textbf{Fig. 4.}\ Clustering and boundary extraction results after
the exploration phase. The original shapes (in blue) are shown
alongside the estimated ones (in red).

\vspace{0.8em}

% ---------- Section IV-D continuation ----------
shown in Fig.~4. A total of 50 simulations were conducted under
identical conditions, incorporating additive white noise
$\mathcal{N}(0, 6.25\;\mathrm{kPa}^{2})$. Each run was terminated
after 70 s to ensure consistent comparison across trials. For each
run, the estimated stiffness map was segmented using the clustering
and boundary extraction pipeline described earlier. The resulting
segmented areas were compared with ground truth using two standard
metrics [15], sensitivity and specificity, defined as:
\[
\text{1) Sensitivity} = \frac{\text{TP}}{\text{TP}+\text{FN}},
\quad
\text{2) Specificity} = \frac{\text{TN}}{\text{TN}+\text{FP}},
\]
where TP, FP, FN, and TN denote the number of pixels of true
positives, false positives, false negatives, and true negatives,
respectively. The metrics were computed separately for each of the
three ground-truth regions, and the final values were averaged
across the simulations. The results of the segmentation algorithm
(Tab.~III) demonstrate elevated and consistent specificity across
all three clusters for the two methods that sample continuously,
thereby confirming the model effectiveness in correctly excluding
background regions. Conversely, the ergodic search exhibited
superior sensitivity compared to both the BO with and without
sampling when clustering non-uniform shapes (Clusters 1--3). This
finding indicates that the ergodic search is more adept at
accurately delineating regions of interest. This is especially
relevant in medical contexts, where false negatives can lead to
missed detections of pathological regions. Consequently, achieving
high sensitivity is imperative to ensure clinically reliable
performance.

% ---------- Section V ----------
\section*{V.\ Experimental Results}

The validity of the simulation results in Sec.~IV was assessed in a
realistic setting by testing the proposed method on a manipulator
interacting with a physical sample (Fig.~1). An UR3e, position-%
controlled at 500 Hz, was used to execute the trajectories generated
by the Ergodic Control. The force at the end-effector was measured
by a 6-axis F/T sensor, a Bota SensOne, sampling at 1 kHz. A 3D-%
printed indenter with a spherical tip with 5 mm radius is attached
to the end-effector. The communication between the robot, the F/T
sensor, and the algorithms was implemented with ROS2 Humble. To
simulate a biological tissue with a mass, we fabricated a silicone
sample consisting of a soft matrix made from Ecoflex-0030 and a
stiffer spherical inclusion (radius 1 cm) made from DragonSkin-30NV.
This configuration increases local stiffness by up to 90\%,
consistent with values reported in previous studies [15], [29]. The
ground truth is built with point-wise sampling on a grid with a
palpation every 2.5 mm in a square of 5 cm $\times$ 5 cm.

We conducted 5 experiments on the physical setup, each starting from
a different random initial position. In all cases, the algorithm
successfully identified the location of the spherical intrusion,
demonstrating consistent target localisation. The reconstruction
accuracy (RMSE $= 9.76 \pm 2.97$ kPa) shows good agreement with the
reference stiffness distribution obtained from grid-based probing
used as ground truth. The method also achieved good classification
performance, with an average sensitivity of 0.92, and specificity of
0.98. Compared to simulation results, the experimental metrics
confirm the high sensitivity and specificity, which are particularly
important for avoiding false negatives in medical applications.
Despite the experimental RMSE values being marginally higher than
those observed in simulation, the measurement uncertainty is still
below the 5\% (average absolute error $8.00 \pm 2.7$ kPa) and the
clustering performance appears largely unaffected. This difference
can be attributed to several factors, including the non-ideal
behaviour of real materials, noise in force measurements, slight
deviations from the assumed contact model, and uncertainties in
ground-truth stiffness calibration. Despite these challenges, the
observed trends are consistent with the simulation, confirming the
robustness of the method in real-world conditions. As shown in
Fig.~5, the trajectory of the end-effector adapts to the underlying
stiffness distribution: it concentrates more exploration in the
stiffer region on the left while still covering the rest of the map.

\end{document}
